\subsection{Estrutura do controle PD}
Seja o erro $e(t)=r(t)-y(t)$ entre a referência $r(t)$ e a saída $y(t)$.
A lei de controle proporcional derivativa é escrita como:
\begin{equation}
\label{eq:pd_law}
u(t) = K_p e(t) + K_d \dot e(t),
\end{equation}
em que $K_p$ e $K_d$ são ganhos a serem ajustados para atender ao comportamento desejado de cada malha.

\subsection{Controle PD no movimento orbital relativo}
No referencial LVLH, o movimento relativo linearizado é descrito pelas equações de HCW.  
Nas aproximações ao longo do eixo radial e para curtas separações, o controlador PD é aplicado em cada eixo do movimento relativo do chaser em relação ao target.

\subsection{Controle PD e a atitude do chaser}
Na malha de atitude, a referência é a orientação desejada para o \emph{chaser}, a saída é a orientação estimada a partir dos sensores de orientação.
O erro de atitude é definido pela diferença entre essas duas informações, sendo representados em quaternions.
% Neste trabalho foi utilizado quaternion para correção do erro entre as orientações.
% O controle PD atua diretamente sobre esse erro: o termo proporcional busca corrigir o desvio de orientação, enquanto o termo derivativo utiliza a velocidade angular ($\vec \omega$) calculada para corrigir a orientação e evitar oscilações excessivas.
O sinal resultante é convertido em torque de controle ($\vec \tau$) e aplicado nos atuadores do \emph{chaser}.
% A sintonia dos ganhos é realizada através de simulação numérica, ajustando $K_p$ e $K_d$ de forma iterativa até se obter uma resposta que atenda aos critérios de tempo de acomodação, estabilidade e esforço de controle.
% Os detalhes completos da dinâmica de atitude serão apresentados em seção posterior, sendo aqui destacada apenas a forma de aplicação do controlador.

\subsection{PD e as juntas do manipulador robótico}
No caso das juntas do manipulador robótico, a referência corresponde ao ângulo desejado de cada junta, enquanto a saída é a posição ou deslocamento medido.
O erro é definido como a diferença entre essas duas variáveis, e a sua derivada representa a velocidade da junta.
O controle PD é aplicado diretamente a esse erro, de modo que o termo proporcional reduz o desvio de posição e o termo derivativo fornece amortecimento à resposta, reduzindo a quantidade de oscilações.
O sinal de controle gerado é aplicado como torque ($\tau_i$) nas juntas revolutas ou como força ($F_{d5}$) na junta prismática.
% O controle PD foi projetado para respeitar os limites de saturação e restrições físicas para que a simulação numérica não se torne irreal.

% Assim como no caso da atitude, a formulação detalhada da dinâmica das juntas será apresentada em seção específica, enquanto que aqui, registra-se a aplicação prática do PD como método de controle.